\documentclass[11pt]{article}
\usepackage[T1]{fontenc}
\usepackage[utf8]{inputenc}
\usepackage{lmodern}
\usepackage{microtype}
\usepackage[margin=0.72in,headheight=15pt]{geometry}
\usepackage[table]{xcolor}
\usepackage{array,booktabs,longtable,tabularx}
\usepackage{amssymb}
\usepackage{enumitem}
\usepackage[most]{tcolorbox}
\usepackage{fancyhdr}
\usepackage{titlesec}
\usepackage{needspace}
\usepackage{ragged2e}
\usepackage{hyperref}
\usepackage{bookmark}
\usepackage{lastpage}

\definecolor{Navy}{HTML}{17324D}
\definecolor{Teal}{HTML}{157A7A}
\definecolor{CueGray}{HTML}{EEF2F4}
\definecolor{SoftBlue}{HTML}{E9F3F8}
\definecolor{SoftGreen}{HTML}{EAF5F0}
\definecolor{SoftAmber}{HTML}{FFF4D6}
\definecolor{RuleGray}{HTML}{B7C3CA}
\definecolor{TextGray}{HTML}{263238}

\hypersetup{colorlinks=true,linkcolor=Navy,urlcolor=Teal,citecolor=Navy,pdfborder={0 0 0}}
\color{TextGray}
\setlength{\parindent}{0pt}
\setlength{\parskip}{4pt}
\setlength{\emergencystretch}{3em}
\hfuzz=0.2pt
\hbadness=2000
\setlist{nosep,leftmargin=1.4em}
\setlength{\LTpre}{4pt}
\setlength{\LTpost}{6pt}
\renewcommand{\arraystretch}{1.2}

\titleformat{\section}{\Large\bfseries\color{Navy}\RaggedRight}{\thesection}{0.55em}{}
\titleformat{\subsection}{\normalsize\bfseries\color{Teal}\RaggedRight}{\thesubsection}{0.5em}{}
\titlespacing*{\section}{0pt}{1.3ex plus .3ex}{.7ex}
\titlespacing*{\subsection}{0pt}{1.1ex plus .2ex}{.5ex}

\pagestyle{fancy}
\fancyhf{}
\lhead{\small\color{Navy}\textbf{Cybersecurity Cornell Notes}}
\rhead{\small\color{Teal}\ChapterMark}
\cfoot{\small\thepage\ of \pageref*{LastPage}}
\renewcommand{\headrulewidth}{0.5pt}

\newtcolorbox{orientationbox}{enhanced,breakable,colback=SoftBlue,colframe=Teal,boxrule=.7pt,arc=2mm,left=2mm,right=2mm,top=1.5mm,bottom=1.5mm}
\newtcolorbox{topicsummary}[1][]{enhanced,breakable,colback=CueGray,colframe=RuleGray,title={Topic Summary},fonttitle=\bfseries\color{Navy},boxrule=.5pt,arc=1.5mm,left=2mm,right=2mm,top=1mm,bottom=1mm,#1}
\newtcolorbox{defensebox}{enhanced,breakable,colback=SoftGreen,colframe=Teal,title={Defensive Guidance},fonttitle=\bfseries,boxrule=.7pt,arc=2mm}
\newtcolorbox{historybox}{enhanced,breakable,colback=SoftAmber,colframe=orange!70!black,title={Historical Context},fonttitle=\bfseries,boxrule=.7pt,arc=2mm}
\newtcolorbox{synthesisbox}{enhanced,breakable,colback=Navy!5,colframe=Navy,title={Chapter Synthesis},fonttitle=\bfseries,boxrule=.8pt,arc=2mm}
\newtcolorbox{takeawaybox}{enhanced,colback=Teal!8,colframe=Teal,title={One-Sentence Takeaway},fonttitle=\bfseries,boxrule=.9pt,arc=2mm}

\newcommand{\CornellHeader}{%
\rowcolor{Navy}\color{white}\bfseries Cue / Recall Prompt & \color{white}\bfseries Notes / Analysis\\\hline}
\newcommand{\CornellRow}[2]{%
\cellcolor{CueGray}\RaggedRight\textbf{#1} & \RaggedRight #2\\\hline}

\newcommand{\ChapterMark}{Chapter 73}
\title{\textcolor{Navy}{Chapter 73: Online Identity and User Management Services}\\[2mm]\large Cornell Notes}
\author{}
\date{}
\hypersetup{pdftitle={Chapter 73: Online Identity and User Management Services - Cornell Notes},pdfauthor={},pdfsubject={Cybersecurity study notes},pdfkeywords={identity, user, management, and mobile}}
\begin{document}
\maketitle
\thispagestyle{fancy}
\vspace{-1.5em}
\begin{orientationbox}
\textbf{Chapter orientation.} This chapter examines online identity and user management services within the broader domain of practical security. Its central problem is how to reason about identity, user, management, and mobile while keeping security objectives, operating assumptions, evidence, and recovery connected. A practitioner should approach the material through decision rights, accountable ownership, risk treatment, policy enforcement, measurement, and assurance. Useful prerequisites include basic risk, threat, control, and systems concepts.
\end{orientationbox}
\section*{Learning Objectives}
\begin{itemize}
\item Explain the security purpose and scope of Online Identity and User Management Services.
\item Identify the assets, actors, assumptions, and trust boundaries associated with identity, user, and management.
\item Distinguish Introduction from Evolution Of Identity Management Requirements and explain where their responsibilities intersect.
\item Analyze how failures in identity, user, and management can affect confidentiality, integrity, availability, privacy, or safety.
\item Evaluate relevant controls through decision rights, accountable ownership, risk treatment, policy enforcement, measurement, and assurance.
\item Audit an implementation using approved policies, ownership records, risk decisions, control tests, exceptions, metrics, and management review.
\item Prioritize remediation according to exploitability, consequence, control strength, and recovery readiness.
\item Verify that corrective actions change observable risk rather than documentation alone.
\end{itemize}
\section*{Cornell Cue-and-Notes}
\Needspace{8\baselineskip}
\subsection*{Introduction}
\begin{longtable}{>{\RaggedRight\arraybackslash}p{0.29\textwidth}|>{\RaggedRight\arraybackslash}p{0.65\textwidth}}
\CornellHeader
\CornellRow{What security problem does Introduction address?}{\textbf{Core idea.} Treat Introduction as a structured part of Online Identity and User Management Services, with particular attention to world, ubiquitous computing, management ambient, and intelligent ubiquitous. \textbf{Analytical lens.} This topic establishes a component of the chapter's argument; connect its definitions and mechanisms to a testable security objective. For world, ubiquitous computing, and management ambient, require evidence that policy intent changes decisions and operating behavior. \textbf{Security reasoning.} Identify what must remain protected, who or what can alter the expected state, which assumptions cross a trust boundary, and how failure changes confidentiality, integrity, availability, privacy, or safety. \textbf{Application.} The practitioner should reconcile intent, implementation, evidence, exceptions, and accountable approval. \textbf{Evidence.} Seek approved policies, ownership records, risk decisions, control tests, exceptions, metrics, and management review.}
\end{longtable}
\begin{topicsummary}
Introduction is best remembered as the connection between world, ubiquitous computing, management ambient, and intelligent ubiquitous and the chapter's larger control objective. A defensible review states the assumption, expected behavior, evidence, failure consequence, and required response.
\end{topicsummary}
\Needspace{8\baselineskip}
\subsection*{Evolution Of Identity Management Requirements}
\begin{longtable}{>{\RaggedRight\arraybackslash}p{0.29\textwidth}|>{\RaggedRight\arraybackslash}p{0.65\textwidth}}
\CornellHeader
\CornellRow{Which assumptions matter in Evolution Of Identity Management Requirements?}{\textbf{Core idea.} Treat Evolution Of Identity Management Requirements as a structured part of Online Identity and User Management Services, with particular attention to identity, management, user, and digital. \textbf{Analytical lens.} This topic is identity-centered: distinguish identification, authentication, authorization, session state, privilege, and accountability. For identity, management, and user, require evidence that policy intent changes decisions and operating behavior. \textbf{Security reasoning.} Identify what must remain protected, who or what can alter the expected state, which assumptions cross a trust boundary, and how failure changes confidentiality, integrity, availability, privacy, or safety. \textbf{Application.} The practitioner should reconcile intent, implementation, evidence, exceptions, and accountable approval. \textbf{Evidence.} Seek approved policies, ownership records, risk decisions, control tests, exceptions, metrics, and management review. Related subtopics include Digital Identity Definition, Identity Management Overview, Privacy Requirement, and User Centricity.}
\end{longtable}
\begin{topicsummary}
Evolution Of Identity Management Requirements is best remembered as the connection between identity, management, user, and digital and the chapter's larger control objective. A defensible review states the assumption, expected behavior, evidence, failure consequence, and required response.
\end{topicsummary}
\Needspace{8\baselineskip}
\subsection*{Identity Management 1.0}
\begin{longtable}{>{\RaggedRight\arraybackslash}p{0.29\textwidth}|>{\RaggedRight\arraybackslash}p{0.65\textwidth}}
\CornellHeader
\CornellRow{How should a reviewer evaluate Identity Management 1.0?}{\textbf{Core idea.} Treat Identity Management 1.0 as a structured part of Online Identity and User Management Services, with particular attention to password, identity, authentication, and user. \textbf{Analytical lens.} This topic is identity-centered: distinguish identification, authentication, authorization, session state, privilege, and accountability. For password, identity, and authentication, require evidence that policy intent changes decisions and operating behavior. \textbf{Security reasoning.} Identify what must remain protected, who or what can alter the expected state, which assumptions cross a trust boundary, and how failure changes confidentiality, integrity, availability, privacy, or safety. \textbf{Application.} The practitioner should reconcile intent, implementation, evidence, exceptions, and accountable approval. \textbf{Evidence.} Seek approved policies, ownership records, risk decisions, control tests, exceptions, metrics, and management review.}
\end{longtable}
\begin{topicsummary}
Identity Management 1.0 is best remembered as the connection between password, identity, authentication, and user and the chapter's larger control objective. A defensible review states the assumption, expected behavior, evidence, failure consequence, and required response.
\end{topicsummary}
\Needspace{8\baselineskip}
\subsection*{The Requirements Fulfilled By Identity Management Technologies}
\begin{longtable}{>{\RaggedRight\arraybackslash}p{0.29\textwidth}|>{\RaggedRight\arraybackslash}p{0.65\textwidth}}
\CornellHeader
\CornellRow{Why can The Requirements Fulfilled By Identity Management Technologies fail in practice?}{\textbf{Core idea.} Treat The Requirements Fulfilled By Identity Management Technologies as a structured part of Online Identity and User Management Services, with particular attention to identity, user, management, and federated. \textbf{Analytical lens.} This topic is identity-centered: distinguish identification, authentication, authorization, session state, privilege, and accountability. For identity, user, and management, require evidence that policy intent changes decisions and operating behavior. \textbf{Security reasoning.} Identify what must remain protected, who or what can alter the expected state, which assumptions cross a trust boundary, and how failure changes confidentiality, integrity, availability, privacy, or safety. \textbf{Application.} The practitioner should reconcile intent, implementation, evidence, exceptions, and accountable approval. \textbf{Evidence.} Seek approved policies, ownership records, risk decisions, control tests, exceptions, metrics, and management review. Related subtopics include Solution by Aggregation, Silo Model, Centralized Versus Federation Identity, and Management.}
\end{longtable}
\begin{topicsummary}
The Requirements Fulfilled By Identity Management Technologies is best remembered as the connection between identity, user, management, and federated and the chapter's larger control objective. A defensible review states the assumption, expected behavior, evidence, failure consequence, and required response.
\end{topicsummary}
\Needspace{8\baselineskip}
\subsection*{Social Login And User Management}
\begin{longtable}{>{\RaggedRight\arraybackslash}p{0.29\textwidth}|>{\RaggedRight\arraybackslash}p{0.65\textwidth}}
\CornellHeader
\CornellRow{What security problem does Social Login And User Management address?}{\textbf{Core idea.} Treat Social Login And User Management as a structured part of Online Identity and User Management Services, with particular attention to higgins, identity, user, and management. \textbf{Analytical lens.} This topic is governance-centered: connect required intent to accountable ownership, implementation, evidence, exceptions, and corrective action. For higgins, identity, and user, require evidence that policy intent changes decisions and operating behavior. \textbf{Security reasoning.} Identify what must remain protected, who or what can alter the expected state, which assumptions cross a trust boundary, and how failure changes confidentiality, integrity, availability, privacy, or safety. \textbf{Application.} The practitioner should reconcile intent, implementation, evidence, exceptions, and accountable approval. \textbf{Evidence.} Seek approved policies, ownership records, risk decisions, control tests, exceptions, metrics, and management review.}
\end{longtable}
\begin{topicsummary}
Social Login And User Management is best remembered as the connection between higgins, identity, user, and management and the chapter's larger control objective. A defensible review states the assumption, expected behavior, evidence, failure consequence, and required response.
\end{topicsummary}
\Needspace{8\baselineskip}
\subsection*{Identity 2.0 For Mobile Users}
\begin{longtable}{>{\RaggedRight\arraybackslash}p{0.29\textwidth}|>{\RaggedRight\arraybackslash}p{0.65\textwidth}}
\CornellHeader
\CornellRow{Which assumptions matter in Identity 2.0 For Mobile Users?}{\textbf{Core idea.} Treat Identity 2.0 For Mobile Users as a structured part of Online Identity and User Management Services, with particular attention to mobile, identity, user, and management. \textbf{Analytical lens.} This topic is identity-centered: distinguish identification, authentication, authorization, session state, privilege, and accountability. For mobile, identity, and user, require evidence that policy intent changes decisions and operating behavior. \textbf{Security reasoning.} Identify what must remain protected, who or what can alter the expected state, which assumptions cross a trust boundary, and how failure changes confidentiality, integrity, availability, privacy, or safety. \textbf{Application.} The practitioner should reconcile intent, implementation, evidence, exceptions, and accountable approval. \textbf{Evidence.} Seek approved policies, ownership records, risk decisions, control tests, exceptions, metrics, and management review. Related subtopics include Introduction, Mobile Web 2.0, Mobility, and Evolution of Mobile Identity.}
\end{longtable}
\begin{topicsummary}
Identity 2.0 For Mobile Users is best remembered as the connection between mobile, identity, user, and management and the chapter's larger control objective. A defensible review states the assumption, expected behavior, evidence, failure consequence, and required response.
\end{topicsummary}
\begin{historybox}
Some products, protocols, examples, threat observations, or legal references in this chapter are historically bounded. Retain the underlying threat model and control objective, but verify present applicability before treating a named implementation as current guidance.
\end{historybox}
\begin{defensebox}
Apply the chapter by making assets, authority, trust, expected behavior, and failure response explicit. In practice, reconcile intent, implementation, evidence, exceptions, and accountable approval. A control is not fully supported until its operation, monitoring, ownership, and recovery behavior are demonstrated with evidence.
\end{defensebox}
\section*{Threat-and-Control Map}
\begingroup\scriptsize\setlength{\tabcolsep}{2pt}
\begin{longtable}{>{\RaggedRight\arraybackslash}p{0.135\textwidth}>{\RaggedRight\arraybackslash}p{0.145\textwidth}>{\RaggedRight\arraybackslash}p{0.155\textwidth}>{\RaggedRight\arraybackslash}p{0.145\textwidth}>{\RaggedRight\arraybackslash}p{0.16\textwidth}>{\RaggedRight\arraybackslash}p{0.17\textwidth}}
\toprule\rowcolor{Navy}\color{white}\textbf{Asset} & \color{white}\textbf{Threat / Failure} & \color{white}\textbf{Vulnerability} & \color{white}\textbf{Impact} & \color{white}\textbf{Detection / Evidence} & \color{white}\textbf{Control}\\\midrule
\endfirsthead
\toprule\rowcolor{Navy}\color{white}\textbf{Asset} & \color{white}\textbf{Threat / Failure} & \color{white}\textbf{Vulnerability} & \color{white}\textbf{Impact} & \color{white}\textbf{Detection / Evidence} & \color{white}\textbf{Control}\\\midrule
\endhead
Governance decisions and organizational assurance & Unmanaged or inconsistently treated risk & Unclear ownership, incomplete policy, or weak measurement & Control gaps, poor decisions, and avoidable exposure & Audit evidence, metrics, exceptions, and management review & Defined authority, risk-based policy, testing, and corrective action\\\midrule
Assumptions surrounding identity & Misuse or unexpected state transition & Trust in user is not validated & A local weakness becomes a broader compromise path & Review evidence for management and test negative paths & Make trust explicit, constrain authority, and fail safely\\\midrule
Recovery capability and decision evidence & Control failure remains unnoticed or cannot be reversed & Monitoring, ownership, or restoration has not been exercised & Longer exposure and slower, less reliable recovery & Alert-to-response exercises and restoration tests & Assigned ownership, actionable telemetry, preserved evidence, and rehearsed recovery\\\midrule
\bottomrule
\end{longtable}
\endgroup
\section*{Practical Security Checklist}
\begin{itemize}[label=$\square$]
\item Define the in-scope assets, users, services, data, and operational dependencies for Online Identity and User Management Services.
\item Document the expected behavior and security objective of identity.
\item Map identities, privileges, trust boundaries, and externally controlled inputs associated with identity and user.
\item Record the threat or failure being addressed, its prerequisites, likely path, and credible consequences.
\item Inspect approved policies, ownership records, risk decisions, control tests, exceptions, metrics, and management review.
\item Test both the intended path and negative paths, including invalid state, partial failure, excessive privilege, and unavailable dependencies.
\item Confirm preventive controls are complemented by actionable detection and a named response owner.
\item Exercise containment, restoration, and management; record actual recovery behavior and unresolved dependencies.
\item Validate exceptions, compensating controls, expiration dates, and accountable approval.
\item Preserve reproducible evidence and tie each finding to affected assets, risk, remediation, and verification criteria.
\end{itemize}
\begin{synthesisbox}
The chapter's principal lesson is that online identity and user management services must be evaluated as an operating security system rather than an isolated product or statement. The most important failure patterns involve identity, user, management, and mobile, especially when ownership, boundaries, telemetry, or recovery remain implicit. A reviewer should connect architecture and behavior to approved policies, ownership records, risk decisions, control tests, exceptions, metrics, and management review, prioritize findings by reachable consequence and control independence, and verify remediation through observable change. Within Practical Security, this chapter contributes a reusable method: state the objective, expose assumptions, test failure, preserve evidence, and learn from operation.
\end{synthesisbox}
\section*{Self-Test Questions}
\begin{enumerate}
\item What is the central security objective of Online Identity and User Management Services?
\item How does Introduction contribute to the chapter's broader security argument?
\item Distinguish Evolution Of Identity Management Requirements from Identity Management 1.0.
\item Which trust assumptions should be made explicit when evaluating identity?
\item How could a weakness involving user affect confidentiality, integrity, availability, privacy, or safety?
\item What evidence would demonstrate that controls surrounding management operate as intended?
\item Why should prevention, detection, response, and recovery be evaluated together in this domain?
\item Scenario: A business unit follows an informal practice that conflicts with the approved control framework. What should the reviewer examine first, and why?
\item Scenario: A control associated with mobile passes a configuration check but fails during an operational exercise. How should the finding be interpreted and prioritized?
\item How should a practitioner distinguish enduring principles in this chapter from time-sensitive tools, examples, or implementation details?
\end{enumerate}
\Needspace{8\baselineskip}
\section*{Answer Key}
\begin{enumerate}
\item The objective is to protect the relevant assets by understanding decision rights, accountable ownership, risk treatment, policy enforcement, measurement, and assurance and by connecting controls to measurable risk reduction.
\item Introduction provides one part of the causal chain: it should identify expected behavior, dependencies, failure modes, and the evidence needed to justify confidence.
\item Evolution Of Identity Management Requirements and Identity Management 1.0 address different portions of the problem. A strong comparison states their scope, assumptions, enforcement points, evidence, and how a failure in one affects the other.
\item Reviewers should identify who controls identity, what inputs or dependencies it trusts, where privilege changes, what happens during partial failure, and how those assumptions are tested.
\item A weakness in user can propagate across trust boundaries. The actual impact depends on reachable assets, granted authority, control independence, visibility, and recovery capability.
\item Useful evidence includes approved policies, ownership records, risk decisions, control tests, exceptions, metrics, and management review; configuration alone is insufficient when operation, monitoring, or recovery is part of the claim.
\item Prevention reduces probability, detection limits dwell time, response limits spread, and recovery restores objectives. Omitting one layer creates an unexamined failure mode.
\item Begin by establishing scope, ownership, expected behavior, and the violated assumption. Then reconcile intent, implementation, evidence, exceptions, and accountable approval, preserving evidence for prioritization and follow-up.
\item Treat the exercise as stronger evidence than a static check. Prioritize according to reachable impact, likelihood of real operating conditions, missing compensating controls, and recovery difficulty.
\item Retain the principle, threat model, and control objective; separately label technologies, legal references, product examples, and threat observations whose applicability depends on time or environment.
\end{enumerate}
\section*{Key-Term Recap}
\begin{longtable}{>{\RaggedRight\arraybackslash}p{0.27\textwidth}>{\RaggedRight\arraybackslash}p{0.66\textwidth}}
\toprule\rowcolor{Navy}\color{white}\textbf{Term} & \color{white}\textbf{Meaning}\\\midrule
\endfirsthead
\toprule\rowcolor{Navy}\color{white}\textbf{Term} & \color{white}\textbf{Meaning}\\\midrule
\endhead
\textbf{Identity Management} & The lifecycle processes and technologies for identities, credentials, attributes, entitlements, and access decisions. \\\midrule
\textbf{Authentication} & The process of establishing confidence that an asserted identity is genuine. \\\midrule
\textbf{Privacy} & The appropriate governance of information about people, including collection, use, disclosure, retention, and individual interests. \\\midrule
\textbf{Risk} & The effect of uncertainty on objectives, commonly evaluated through likelihood, impact, exposure, and context. \\\midrule
\textbf{Policy} & A management statement of required intent, responsibilities, and acceptable behavior. \\\midrule
\textbf{Authorization} & The decision process that determines which authenticated subject may perform which action on a resource. \\\midrule
\textbf{Certificate} & A signed data structure that binds identity information to a public key under a trust model. \\\midrule
\textbf{Biometrics} & Identity evidence derived from physiological or behavioral characteristics, evaluated probabilistically rather than as a secret. \\\midrule
\textbf{Access Control} & Rules and enforcement mechanisms that restrict actions on resources to authorized subjects. \\\midrule
\textbf{Integrity} & The property that information and systems remain accurate, complete, and protected from unauthorized modification. \\\midrule
\bottomrule
\end{longtable}
\begin{takeawaybox}
Treat online identity and user management services as a verifiable chain from assets and assumptions through controls, evidence, response, and recovery.
\end{takeawaybox}
\end{document}
